\documentclass[11pt]{article}

\usepackage{amsmath,amssymb,amsthm}
\usepackage{mathtools}
\usepackage[margin=1in]{geometry}
\usepackage{booktabs}
\usepackage{hyperref}

\newtheorem{lemma}{Lemma}
\newtheorem{remark}{Remark}

\title{From Sum-Rate to Energy-Efficient Beamforming:\\
       What Is Reused, What Is Extended, and What Is Novel}
\author{}
\date{}

\begin{document}
\maketitle

\section{Purpose of this note}

This note separates three things that are easy to blur together when
presenting the idea: (i) what is taken directly from the reference paper
\cite{Schroder2025}, (ii) what we changed, and (iii) why the change is a
genuine departure and not a relabeling of the same optimum. The system and
channel model in Section~\ref{sec:system} is unchanged from the paper and is
restated only to fix notation. Section~\ref{sec:related} positions this
extension against the closest related work found outside the reference
paper's own line of work.

\section{System model (adopted, unchanged)}
\label{sec:system}

A set $\mathcal{M}=\{1,\dots,M\}$ of LEO satellites, each equipped with a
uniform linear array of $N$ antennas, serves a set $\mathcal{K}=\{1,\dots,K\}$
single-antenna users in the downlink. The channel between satellite $m$ and
user $k$ is $\mathbf{h}_{k,m}\in\mathbb{C}^{1\times N}$, stacked as
$\mathbf{h}_k = [\mathbf{h}_{k,1},\dots,\mathbf{h}_{k,M}] \in
\mathbb{C}^{1\times MN}$. Data symbols are precoded with
$\mathbf{w}_k \in \mathbb{C}^{MN\times 1}$ and superimposed,
$\mathbf{x} = \sum_{k\in\mathcal{K}} \mathbf{w}_k u_k$. The full precoding
matrix is $\mathbf{W} = [\mathbf{w}_1,\dots,\mathbf{w}_K] \in
\mathbb{C}^{MN\times K}$. The instantaneous achievable sum rate is

\begin{equation}
  R(\mathbf{W}) \;=\; \sum_{k\in\mathcal{K}}
  \log_2\!\left(1 + \frac{|\mathbf{h}_k \mathbf{w}_k|^2}
  {\sigma_{\text{noise}}^2 + \sum_{k'\neq k}|\mathbf{h}_k \mathbf{w}_{k'}|^2}\right).
  \label{eq:sumrate}
\end{equation}

CSIT is corrupted by two error sources acting on the true channel
$\mathbf{h}_{k,m}$: an angle-of-departure (steering) error
$\boldsymbol{\epsilon}_{\text{aod},k,m}\sim\mathcal{U}(-\Delta\epsilon_{\text{aod}},
\Delta\epsilon_{\text{aod}})$ and a synchronization phase error
$\epsilon_{\text{ph},k,m}\sim\mathcal{N}(0,\sigma_{\text{ph}}^2)$, giving an
estimate $\tilde{\mathbf{h}}_{k,m}$ used by every precoder in place of the
true $\mathbf{h}_{k,m}$.

\section{The paper's optimization problem (reproduced)}
\label{sec:paper-problem}

The reference paper poses a single objective: maximize expected sum rate
under a per-satellite radiated-power budget,

\begin{equation}
  \max_{\mathbf{W}} \;\; \bar{R} = \mathbb{E}_{\mathbf{H}}[R(\mathbf{W})]
  \qquad \text{s.t.} \quad \|\mathbf{W}_m\|_F^2 \le P/M \quad \forall m\in\mathcal{M}.
  \label{eq:paper-problem}
\end{equation}

Every precoder evaluated in the paper --- MMSE, robust SLNR, and the
learned SAC policy --- enforces this constraint by \emph{unconditional
rescaling} of the raw precoder to the boundary of the feasible ball:

\begin{equation}
  \mathbf{W}_{\text{paper}} \;=\; \sqrt{\frac{P/M}{\|\mathbf{W}_{\text{raw}}\|_F^2}}\;\mathbf{W}_{\text{raw}},
  \label{eq:paper-normalization}
\end{equation}

independent of $\|\mathbf{W}_{\text{raw}}\|_F^2$. This is Eq.~(26)/(30) of the
paper for MMSE, and the SAC output transform is normalized identically. The
reward signal driving SAC's actor is the sum rate $R(\mathbf{W}_{\text{paper}})$
alone; power never appears as a cost, only as this hard, always-active
rescaling.

\begin{lemma}[Full-power optimality under rate-only objectives]
\label{lem:saturate}
Fix a precoding \emph{direction} $\hat{\mathbf{W}}$ with
$\|\hat{\mathbf{W}}\|_F^2=1$ and consider the scaled family
$\mathbf{W}(\alpha) = \alpha\hat{\mathbf{W}}$, $\alpha\in(0,\sqrt{P/M}]$. Then
$R(\mathbf{W}(\alpha))$ is non-decreasing in $\alpha$, so the maximizer of
\eqref{eq:paper-problem} over this family is $\alpha^\star = \sqrt{P/M}$.
\end{lemma}

\begin{proof}[Sketch]
Scaling $\mathbf{W}$ by $\alpha$ scales every signal term
$|\mathbf{h}_k\mathbf{w}_k|^2$ and every interference term
$|\mathbf{h}_k\mathbf{w}_{k'}|^2$ by $\alpha^2$, while
$\sigma_{\text{noise}}^2$ is fixed. Hence each user's SINR,
$\alpha^2|\mathbf{h}_k\hat{\mathbf{w}}_k|^2 /
(\sigma_{\text{noise}}^2 + \alpha^2\sum_{k'\neq k}|\mathbf{h}_k\hat{\mathbf{w}}_{k'}|^2)$,
is non-decreasing in $\alpha$ (strictly increasing while noise-limited,
saturating to an interference-limited ceiling as $\alpha\to\infty$, never
decreasing). $\log_2(1+\cdot)$ is monotone, so $R$ inherits the same
monotonicity. $\blacksquare$
\end{proof}

Lemma~\ref{lem:saturate} is not a limitation of the paper's method --- it is
the \emph{correct} optimal behavior for problem~\eqref{eq:paper-problem}.
Rescaling every raw output to the boundary (Eq.~\ref{eq:paper-normalization})
is therefore not a modeling simplification; it exactly implements the
optimal solution structure of a rate-only objective. Any precoder solving
\eqref{eq:paper-problem} \emph{should} saturate the power budget.

\section{Our extension: making power a cost, not just a cap}
\label{sec:our-extension}

\subsection{Power consumption model (new)}

The paper's constraint counts only radiated power, $\|\mathbf{W}\|_F^2$. We
add a satellite HPA electrical power model of the same form used in recent
energy-efficient satellite precoding work \cite{HaEtAl2024}, splitting DC
power draw into amplifier draw and static circuit dissipation:

\begin{equation}
  P_{\text{total}}(\mathbf{W}) \;=\; \underbrace{\frac{\mathrm{tr}\{\mathbf{W}^H\mathbf{W}\}}{\eta_{\text{PA}}}}_{\text{amplifier draw}}
  \;+\; \underbrace{MN\,P_c}_{\text{circuit draw}},
  \label{eq:ptotal}
\end{equation}

with power-amplifier efficiency $\eta_{\text{PA}}\in(0,1)$ and per-antenna
static circuit power $P_c$. We set $\eta_{\text{PA}}=0.60$, matching the
satellite antenna/HPA efficiency used in \cite{HaEtAl2024} and consistent
with reported saturated TWTA/SSPA efficiencies of 55--70\%
\cite{Lohmeyer2016}. An earlier version of this model used
$\eta_{\text{PA}}=0.35$ sourced from a \emph{terrestrial}
macro/micro/pico/femto base-station power study \cite{Auer2011} -- not
satellite hardware -- and has been corrected. Two structural properties of
\eqref{eq:ptotal} matter for what follows:

\begin{itemize}
  \item \textbf{No such term exists in the reference paper.} Its optimization
    problem \eqref{eq:paper-problem} has no notion of amplifier efficiency or
    circuit power; $P_{\text{total}}$ is our addition in full.
  \item $P_{\text{total}}$ has a strictly positive floor $MN P_c/\eta_{\text{PA}}\cdot 0 + MNP_c = MNP_c$
    as $\mathrm{tr}\{\mathbf{W}^H\mathbf{W}\}\to 0$: the circuit term never
    vanishes. This bounds the maximum achievable percentage power savings
    and keeps any rate/power ratio well defined even at zero transmit power.
\end{itemize}

\subsection{The energy-efficiency objective (new)}

We replace \eqref{eq:paper-problem} with a fractional (concave-over-affine)
objective,

\begin{equation}
  \max_{\mathbf{W}} \;\; \mathrm{EE}(\mathbf{W}) = \frac{R(\mathbf{W})}{P_{\text{total}}(\mathbf{W})}
  \qquad \text{s.t.} \quad \|\mathbf{W}_m\|_F^2 \le P/M \quad \forall m\in\mathcal{M}.
  \label{eq:ee-problem}
\end{equation}

This is a genuinely different problem from \eqref{eq:paper-problem}, not a
relaxation of it: power now appears in the objective, not only in the
constraint. Lemma~\ref{lem:saturate} no longer applies, because scaling
$\mathbf{W}(\alpha)=\alpha\hat{\mathbf{W}}$ now moves \emph{both} the
numerator and the denominator:

\begin{equation}
  \mathrm{EE}(\alpha) = \frac{R(\alpha\hat{\mathbf{W}})}
  {\alpha^2\|\hat{\mathbf{W}}\|_F^2/\eta_{\text{PA}} + MNP_c}.
  \label{eq:ee-alpha}
\end{equation}

The numerator saturates (Lemma~\ref{lem:saturate}, bounded above by the
interference-limited ceiling); the denominator grows without bound, linearly
in $\alpha^2$. Consequently $\mathrm{EE}(\alpha)$ is not monotone: it
generically increases while noise-limited, then \emph{decreases} once the
marginal rate gained per additional watt falls below the marginal cost of
that watt. The maximizer $\alpha^\star$ of \eqref{eq:ee-alpha} is generically
in the \emph{interior} of $(0,\sqrt{P/M}]$, not at the boundary. This is the
precise sense in which the idea is not a relabeling of the paper's problem:
we did not merely permit $\|\mathbf{W}\|_F^2 < P/M$ (that permission alone is
vacuous under \eqref{eq:paper-problem}, by Lemma~\ref{lem:saturate}); we
changed the objective so the boundary stops being automatically optimal.

\begin{remark}
This is precisely the classical energy-efficiency/spectral-efficiency
trade-off \cite{Chen2011}: under a rate-only objective, more power is always
weakly better and the power cap binds; under a rate-per-power objective, the
optimal power level is an interior stationary point of a concave-over-affine
ratio, characterized by Dinkelbach's transform (Section~\ref{sec:schemes}).
\end{remark}

\subsection{Relaxing the always-saturate precoder (new)}

To let the learned policy actually reach the interior of the feasible ball,
we replace the unconditional rescaling \eqref{eq:paper-normalization} with a
\emph{clip-only} projection,

\begin{equation}
  \mathbf{W}_{\text{ours}} =
  \begin{cases}
    \mathbf{W}_{\text{raw}}, & \|\mathbf{W}_{\text{raw}}\|_F^2 \le P/M, \\[4pt]
    \sqrt{P/M}/\|\mathbf{W}_{\text{raw}}\|_F \cdot \mathbf{W}_{\text{raw}}, & \text{otherwise.}
  \end{cases}
  \label{eq:clip-only}
\end{equation}

Constraint (15b) of the paper is enforced exactly as written (an inequality)
rather than as the de facto equality every paper precoder produces. This is
the actual, literal feasible set of \eqref{eq:ee-problem}; \eqref{eq:paper-normalization}
only ever explores its boundary.

\section{Three formulations implemented and compared}
\label{sec:schemes}

\paragraph{Scheme I --- reproduction of the paper's convention.} Precoder
\eqref{eq:paper-normalization}, reward $R(\mathbf{W})/P_{\text{total}}$. Since
$P_{\text{total}}$ is a constant here (power is always exactly $P/M$), this
reduces to maximizing $R(\mathbf{W})$ alone --- it is the paper's sum-rate
objective in all but name, retained purely as a same-conditions reference
point.

\paragraph{Scheme II --- direct-ratio reward, free power.} Precoder
\eqref{eq:clip-only}, reward $R(\mathbf{W})/P_{\text{total}}(\mathbf{W})$
evaluated directly as a ratio. Numerically fragile as
$P_{\text{total}}\to MNP_c$ from below is bounded (Section~\ref{sec:our-extension}),
but the ratio's gradient can still be poorly scaled for online SGD.

\paragraph{Scheme III --- Dinkelbach-reformulated reward, free power.}
Precoder \eqref{eq:clip-only}; reward is the Dinkelbach subtractive transform
of the fractional program,

\begin{equation}
  \mathrm{reward}_t = R(\mathbf{W}_t) - \lambda_t\, P_{\text{total}}(\mathbf{W}_t),
  \qquad
  \lambda_{t+1} = \frac{\overline{R}_{\text{window}}}{\overline{P}_{\text{total},\text{window}}},
  \label{eq:dinkelbach}
\end{equation}

the classical fixed-point reformulation \cite{Dinkelbach1967}:
$F(\lambda) = \max_{\mathbf{W}} \left[R(\mathbf{W}) - \lambda P_{\text{total}}(\mathbf{W})\right]$,
with $F(\lambda^\star)=0 \iff \lambda^\star = \mathrm{EE}(\mathbf{W}^\star)$.
Framing \eqref{eq:dinkelbach} as a running stochastic-approximation update
(rather than solving each $\lambda$-subproblem to optimality before updating,
as the textbook alternating scheme prescribes) is a deliberate relaxation for
the online RL setting, made explicit here rather than left implicit.

\section{Theoretical grounding for the online Dinkelbach update}
\label{sec:dinkelbach-rigor}

Framing~\eqref{eq:dinkelbach} as a running stochastic-approximation update was
flagged in Section~\ref{sec:schemes} as a deliberate relaxation of the
classical scheme, in which $\lambda$ is updated only after the inner problem
$\max_{\mathbf{W}}[R(\mathbf{W})-\lambda P_{\text{total}}(\mathbf{W})]$ is
solved to its exact global optimum \cite{Dinkelbach1967}. SAC does neither:
its inner solve is approximate, local, and updated concurrently with
$\lambda$ rather than beforehand. This section asks, and answers, whether
that relaxation has any rigorous grounding, without changing the underlying
algorithm.

\subsection{What the closest existing results do and do not cover}

Four adjacent literatures combine a fractional (ratio) objective with a
learned/iterative policy, and none directly covers a concurrently-updated
outer multiplier with function approximation. The classical beamforming
literature \cite{HaEtAl2024,DengEtAl2023} uses Dinkelbach's method with the
inner problem solved to full, certified convergence (convex optimization or
tabular Q-learning re-solved exactly each step) before $\lambda$ is touched.
Risk-sensitive RL results for ratio objectives such as the Sharpe ratio
\cite{SharpeRatioMDP2025} use the same exact block/outer-loop pattern.
Safe/constrained RL does have real convergence proofs for genuinely
\emph{concurrent} primal-dual updates \cite{Paternain2019,DingZhangBasar2020},
but only for a softmax policy over a finite action space --- not for
continuous actions or neural function approximation. The one result that
formally extends Dinkelbach's method to an \emph{inexact} inner solve,
Jin et al.\ \cite{JinTangZhangWang2023}, requires (a) an explicit,
per-iteration approximation-error bound $\epsilon_i$ that (b) shrinks over
iterations, and (c) convexity of the inner value function in $\lambda$. As
implemented, EE\_sac.py satisfies none of the three as stated for a general
MDP. The remainder of this section shows that, for this specific
environment, (c) holds unconditionally, (a) is obtainable from standard
nonconvex stochastic-optimization theory, and (b) is the one genuine gap ---
closable by a training-schedule change, not an architectural one.

\subsection{The environment is a contextual bandit}

\begin{lemma}[Action-independent transitions collapse temporal credit assignment]
\label{lem:bandit}
Let $s_t$ denote the SAC state at training step $t$ and $a_t$ the
(pre-projection) action. If the state-transition kernel satisfies
$s_{t+1}\sim p(\cdot\mid s_t)$, independent of $a_t$, then for every state $s$,
$$
Q^\star(s,a) \;=\; r(s,a) \;+\; \gamma\,\mathbb{E}_{s'\sim p(\cdot\mid s)}\big[V^\star(s')\big],
\qquad
\arg\max_a Q^\star(s,a) \;=\; \arg\max_a r(s,a).
$$
\end{lemma}
\begin{proof}
By the Bellman optimality decomposition,
$Q^\star(s,a)=\mathbb{E}[r_t+\gamma V^\star(s_{t+1})\mid s_t=s,a_t=a]$. Since
$s_{t+1}\sim p(\cdot\mid s_t)$ does not depend on $a_t$, the second term
equals $\gamma\,\mathbb{E}_{s'\sim p(\cdot\mid s)}[V^\star(s')]$ for every
$a$ --- a quantity constant in $a$. Adding an $a$-independent term to $r(s,a)$
does not change its argmax. $\blacksquare$
\end{proof}

\begin{remark}[This holds by construction in this codebase, not by assumption]
\texttt{update\_sim()}, the state-transition routine shared by every training
script in this repository, advances user/satellite positions, rolls CSIT
estimation errors, and recomputes channel state information without ever
receiving the chosen precoder as an argument. The hypothesis of
Lemma~\ref{lem:bandit} is therefore exact, not approximate: whichever
$\mathbf{W}_t$ SAC outputs, the channel realization at step $t+1$ is drawn
from the same exogenous process regardless. There is no delayed consequence
of a precoding decision anywhere in this environment --- it is a per-sample
stochastic optimization problem wearing an RL interface, not a genuine
sequential MDP. This is precisely why a concurrently-updated $\lambda$ is not
automatically ruled out here the way it is in the general MDP results of
Section~\ref{sec:dinkelbach-rigor} above: there is no real temporal credit
assignment to get wrong.
\end{remark}

\subsection{Convexity and smoothness on the constraint set}

\begin{lemma}[Automatic convexity of the Dinkelbach value function]
\label{lem:convex}
Define $F(\lambda) = \max_{\mathbf{W}\in\mathcal{W}}\big[R(\mathbf{W}) -
\lambda P_{\text{total}}(\mathbf{W})\big]$ for the feasible ball
$\mathcal{W}=\{\mathbf{W}:\|\mathbf{W}_m\|_F^2\le P/M\ \forall m\}$. Then $F$
is convex in $\lambda$ on $\mathbb{R}$, regardless of the specific form of
$R$ or $P_{\text{total}}$.
\end{lemma}
\begin{proof}
For each fixed $\mathbf{W}\in\mathcal{W}$, $\lambda\mapsto
R(\mathbf{W})-\lambda P_{\text{total}}(\mathbf{W})$ is affine in $\lambda$.
$F$ is a pointwise supremum of affine functions of $\lambda$, hence convex.
$\blacksquare$
\end{proof}

This discharges condition (c) of \cite{JinTangZhangWang2023} unconditionally.

\begin{lemma}[Smoothness of the per-step Dinkelbach surface]
\label{lem:smooth}
Fix $\lambda\ge0$ and let $f_\lambda(\mathbf{W}) = R(\mathbf{W}) - \lambda
P_{\text{total}}(\mathbf{W})$. On $\mathcal{W}$, $f_\lambda$ is continuously
differentiable with Lipschitz-continuous gradient (finite $L_\lambda$-smooth).
\end{lemma}
\begin{proof}[Sketch]
$\sigma_{\text{noise}}^2>0$ is a fixed positive constant, so every SINR
denominator in \eqref{eq:sumrate} is bounded away from zero uniformly on
$\mathcal{W}$; $R$ is a composition of rational functions (bounded
denominator) with $\log_2(1+\cdot)$, hence real-analytic on $\mathcal{W}$.
$P_{\text{total}}$ \eqref{eq:ptotal} is affine in the quadratic form
$\mathrm{tr}\{\mathbf{W}^H\mathbf{W}\}$, with constant Hessian
$(2/\eta_{\text{PA}})I$. A finite linear combination of functions with
continuous Hessian has bounded (hence Lipschitz) Hessian on the compact set
$\mathcal{W}$. $\blacksquare$
\end{proof}

\begin{remark}[The clip is a projection, not an objective-level discontinuity]
The precoder \eqref{eq:clip-only} is exactly Euclidean projection
$\Pi_{\mathcal{W}}$ onto a convex ball. The composite map
$\mathbf{W}_{\text{raw}}\mapsto f_\lambda(\Pi_{\mathcal{W}}(\mathbf{W}_{\text{raw}}))$,
viewed as a function of the unconstrained actor output, has a discontinuous
Jacobian exactly on the sphere $\|\mathbf{W}_{\text{raw}}\|_F^2=P/M$ ---
$\Pi_{\mathcal{W}}$ itself is what is non-smooth there, not $f_\lambda$. This
is exactly the setting nonconvex \emph{projected}-stochastic-gradient theory
is built for: convergence is stated via the gradient mapping
$G_\lambda(\mathbf{W}) = \big(\mathbf{W}-\Pi_{\mathcal{W}}(\mathbf{W}+\alpha\nabla
f_\lambda(\mathbf{W}))\big)/\alpha$ rather than the raw gradient
\cite{GhadimiLanZhang2016}, sidestepping the projection boundary's
non-smoothness instead of requiring it away. Lemma~\ref{lem:smooth} is
exactly the hypothesis on $f_\lambda$ that framework needs.
\end{remark}

\subsection{The one remaining gap: an unshrinking entropy bias}

Under Lemma~\ref{lem:bandit}, SAC's actor/critic are --- in effect ---
running stochastic projected-gradient ascent directly on $f_\lambda$, so the
gradient-mapping norm after $T$ inner steps admits the standard $O(1/\sqrt{T})$
rate \cite{GhadimiLan2013,GhadimiLanZhang2016}: an explicit, checkable route
to condition (a) of \cite{JinTangZhangWang2023}, in place of an
unquantifiable ``training plateaued.'' What this does \emph{not} yet secure
is condition (b), that the inner-solve error actually shrinks over training
rather than plateauing at a nonzero floor.

SAC's maximum-entropy objective is exactly such a floor. The auto-tuned
temperature $\alpha_{\text{SAC}}$ is adjusted throughout training to hold the
policy's entropy at \texttt{target\_entropy}, fixed at $1.0$ for the entire
run in \texttt{config\_sac\_learner.py} and never annealed. The soft-optimal
policy therefore maximizes $\mathbb{E}[Q(s,a)]+\alpha_{\text{SAC}}\mathcal{H}(\pi(\cdot\mid s))$,
not $\mathbb{E}[Q(s,a)]$ alone, and the resulting gap from the true argmax is
bounded by a term of order $\alpha_{\text{SAC}}$ that does not shrink from
episode $0$ to episode $13{,}000$ under the current schedule --- the one
condition of \cite{JinTangZhangWang2023} not currently met.

\begin{remark}[The resulting recommendation, and why it does not require abandoning SAC]
Annealing \texttt{target\_entropy} (or decaying $\alpha_{\text{SAC}}$) over
training turns this floor into a term that vanishes as training proceeds,
which is exactly what condition (b) requires. This is a schedule change to
the existing training loop, not an architectural one: actor, critic, and
replay buffer are unchanged. That matters because, while the environment is
action-\emph{independent} (Lemma~\ref{lem:bandit}), it is not
\emph{stationary} over the mission timeline --- the satellite's motion drifts
the channel statistics continuously. An offline ``learning-to-optimize''
alternative (a feed-forward network trained once by ordinary backpropagation
against a fixed, i.i.d.-sampled batch of channel realizations) would satisfy
Lemma~\ref{lem:smooth} equally well, but would need an explicit
retraining/adaptation mechanism to track that drift, which SAC's online
training loop provides for free. This is why the object being justified here
is the existing online SAC formulation, not a replacement for it.
\end{remark}

A related, purely empirical stabilization technique is worth noting
separately: the $\lambda$ update in \eqref{eq:dinkelbach} is, mathematically,
a pure integral controller on the rate/power gap, which is known from
control theory to be prone to oscillation and overshoot. PID-Lagrangian
methods from safe RL \cite{StookeAchiamAbbeel2020} add proportional and
derivative terms to exactly this kind of update to damp such oscillation.
This is a complementary, cheap engineering fix for training stability; unlike
the entropy-annealing schedule above, it comes with no convergence proof of
its own and does not by itself close any of the three conditions discussed
above.

\subsection{Summary of the argument}

Lemma~\ref{lem:bandit} reduces the relevant question from general-MDP
convergence --- where every concurrently-updated primal-dual scheme found in
the literature either requires a discrete/softmax policy
\cite{Paternain2019,DingZhangBasar2020} or an exactly-solved inner problem
\cite{HaEtAl2024,DengEtAl2023} --- to bandit-level convergence, which is far
better understood. Lemma~\ref{lem:convex} discharges condition (c) of
\cite{JinTangZhangWang2023} unconditionally. Lemma~\ref{lem:smooth}, combined
with projected-SGD gradient-mapping bounds
\cite{GhadimiLan2013,GhadimiLanZhang2016}, gives an explicit route to
condition (a). What remains open is condition (b), currently blocked only by
a fixed, non-annealed entropy target. Annealing that target is therefore the
single concrete change needed to make the online, reward-embedded Dinkelbach
update \eqref{eq:dinkelbach} a rigorously justified instance of an
established inexact-fractional-programming convergence theorem, rather than
an unjustified heuristic --- without departing from the SAC architecture used
throughout this work. Whether entropy annealing changes the trained policy's
realized $\mathrm{EE}(\mathbf{W})$ in practice is a separate, still-open
empirical question.

\section{Related work}
\label{sec:related}

The reference paper \cite{Schroder2025} itself continues an earlier
single-objective (sum-rate only) line of work by the same group
\cite{Gracla2024}, which introduced the flexible SAC-based precoding
framework this note builds on. Outside that lineage, three adjacent bodies
of work are relevant, and it is worth being precise about what each does and
does not share with the extension proposed here.

\textbf{Energy-efficient beamforming via deep RL.} DRL has been applied
directly to EE-objective beamforming design, e.g.\ uplink beamforming for
long-term EE maximization in cell-free networks via DDPG
\cite{LiNiTianHua2021}. This confirms ``RL for an EE objective'' is an
established combination in general, but not in a satellite downlink setting
and not paired with CSIT-error robustness.

\textbf{Energy-efficient learning for LEO satellites.} DRL has also been
used for EE optimization specifically in LEO satellite systems, e.g.\ joint
beamforming, UAV positioning, and rate-splitting for RSMA-based LEO links
assisted by an active reconfigurable intelligent surface
\cite{YeganehBehroozi2025}, and a risk-aware (CVaR-based) DRL framework for
robust hybrid beamforming in massive-MIMO LEO systems that also targets
energy efficiency \cite{RiskAwareLEO2024}. Both combine EE, robustness, and
LEO satellites as we do, but in materially larger/different system
architectures (RIS- and UAV-assisted, or massive-MIMO hybrid precoding) and
via a different robustness mechanism (risk measures rather than training
under a CSIT error distribution).

\textbf{Fractional programming combined with DRL.} The Dinkelbach
transform's classical, non-learned use in beamforming is developed
comprehensively in \cite{ShenYu2018}. Closer to our combination,
\cite{JiaoFangDing2022} jointly uses fractional programming and DRL for
robust EE beamforming in a WPT-assisted D2D/MISO-NOMA setting; there,
fractional programming and DRL solve separate subproblems in an alternating
scheme, rather than the Dinkelbach multiplier being folded directly into a
single RL reward signal as in \eqref{eq:dinkelbach}.

\textbf{Closest classical (non-learned) analogue.} \cite{HaEtAl2024}
solves an EE precoding problem for a multi-gateway, bent-pipe SATCOM system
using Dinkelbach's method combined with an MMSE-based (WMMSE-style)
transformation and a compressed-sensing relaxation for sparse
feeder-link-beam matching -- the same general tool (Dinkelbach) applied to
a materially different, and more classically tractable, version of the
problem. Three differences separate it from this work: (i) it is a
model-based iterative convex optimization, solved assuming accurate channel
knowledge at every step, not a learned policy; (ii) it does not train
against or evaluate robustness to CSIT \emph{estimation error} -- channel
impairments (phase noise, fading) are simulated, but the algorithm optimizes
assuming it knows the channel it is given; and (iii) Dinkelbach's $\lambda$
is updated via the classical exact alternating scheme (solve the inner QCQP
to convergence, then update $\lambda$), not the online, continuously-updated
reward-embedded form of \eqref{eq:dinkelbach}. This sharpens the gap this
work fills: robustness to erroneous CSIT is difficult to fold into a
closed-form/convex-relaxation solver of this kind, which is precisely why
\cite{Schroder2025} argues for a learned policy in the first place -- an
argument this work extends from the sum-rate objective to the EE one.

None of these combine, on the exact system and CSIT-error-robustness setup
of \cite{Schroder2025}, (i) the electrical power model \eqref{eq:ptotal},
(ii) the resulting EE objective \eqref{eq:ee-problem}, and (iii) an
online, reward-embedded Dinkelbach update \eqref{eq:dinkelbach}. That
combination, on this system, is the specific gap this work fills.

\section{Summary: what is the paper's, what is ours}

\begin{center}
\begin{tabular}{@{}p{0.32\linewidth}p{0.63\linewidth}@{}}
\toprule
\textbf{Component} & \textbf{Status} \\
\midrule
Channel model, LoS geometry, CSIT error sources (Eq.~\ref{eq:sumrate} and
$\S$\ref{sec:system}) & Reproduced unchanged from \cite{Schroder2025}. \\
Sum-rate objective \eqref{eq:paper-problem} and always-saturate precoder
\eqref{eq:paper-normalization} & Reproduced unchanged; retained as Scheme~I,
a same-conditions reference point. \\
SAC algorithm (actor/critic structure, entropy objective, offline replay
buffer) & Reused, same underlying learning algorithm. \\
Electrical power model $P_{\text{total}}$ \eqref{eq:ptotal} (amplifier
efficiency, static circuit power) & \textbf{New.} No analogue in the paper. \\
Energy-efficiency objective \eqref{eq:ee-problem} (rate over total power) &
\textbf{New.} Paper only ever maximizes rate subject to a power cap. \\
Clip-only precoder \eqref{eq:clip-only} realizing the true inequality
constraint & \textbf{New.} Every paper precoder instead implements
\eqref{eq:paper-normalization}. \\
Dinkelbach reward reformulation \eqref{eq:dinkelbach} with an online,
windowed $\lambda$ update & \textbf{New.} Not a fractional program in the
paper at all. \\
Rigorous grounding for the online $\lambda$ update
(Section~\ref{sec:dinkelbach-rigor}) & \textbf{New.} Reduces the question to
a bandit-convergence argument with an automatically-convex, provably-smooth
inner surface; identifies one concrete remaining gap (a non-annealed entropy
target) and its fix, without changing the SAC architecture. \\
\bottomrule
\end{tabular}
\end{center}

The claim of novelty rests on Section~\ref{sec:our-extension}: introducing
$P_{\text{total}}$ into the objective is what breaks
Lemma~\ref{lem:saturate}'s always-saturate conclusion, and Section~\ref{sec:related}
situates that claim against the closest work found elsewhere in the
literature. Whether the trained policy \emph{empirically} finds and exploits
an interior optimum that beats Scheme~I on realized $\mathrm{EE}(\mathbf{W})$
is a separate, still-open empirical question, addressed by evaluation rather
than by the formulation itself. Section~\ref{sec:dinkelbach-rigor} addresses
a different open question --- not whether the trained policy performs well,
but whether the training procedure itself has a defensible theoretical basis
--- and identifies exactly one concrete, not-yet-implemented change
(entropy-target annealing) needed to close it.

\begin{thebibliography}{99}

\bibitem{Schroder2025}
A.~Schr\"oder, S.~Gracla, C.~Bockelmann, D.~W\"ubben, and A.~Dekorsy,
``Model-Free Robust Beamforming in Satellite Downlink Using Reinforcement
Learning,'' \emph{IEEE Open Journal of the Communications Society}, 2025.
arXiv:2510.25393.

\bibitem{Gracla2024}
A.~Schr\"oder, S.~Gracla, M.~R\"oper, D.~W\"ubben, C.~Bockelmann, and
A.~Dekorsy, ``Flexible Robust Beamforming for Multibeam Satellite Downlink
using Reinforcement Learning,'' arXiv:2402.16563, 2024.

\bibitem{Auer2011}
G.~Auer, V.~Giannini, C.~Desset, I.~Godor, P.~Skillermark, M.~Olsson,
M.~A.~Imran, D.~Sabella, M.~J.~Gonzalez, O.~Blume, and A.~Fehske, ``How Much
Energy Is Needed to Run a Wireless Network?,'' \emph{IEEE Wireless
Communications}, vol.~18, no.~5, pp.~40--49, 2011.

\bibitem{Chen2011}
Y.~Chen, S.~Zhang, S.~Xu, and G.~Y.~Li, ``Fundamental Trade-offs on Green
Wireless Networks,'' \emph{IEEE Communications Magazine}, vol.~49, no.~6,
pp.~30--37, 2011.

\bibitem{Dinkelbach1967}
W.~Dinkelbach, ``On Nonlinear Fractional Programming,'' \emph{Management
Science}, vol.~13, no.~7, pp.~492--498, 1967.

\bibitem{ShenYu2018}
K.~Shen and W.~Yu, ``Fractional Programming for Communication
Systems---Part I: Power Control and Beamforming,'' \emph{IEEE Transactions
on Signal Processing}, vol.~66, no.~10, pp.~2616--2630, 2018.

\bibitem{LiNiTianHua2021}
W.~Li, W.~Ni, H.~Tian, and M.~Hua, ``Deep Reinforcement Learning for
Energy-Efficient Beamforming Design in Cell-Free Networks,'' in \emph{Proc.
IEEE Wireless Communications and Networking Conference Workshops
(WCNCW)}, 2021. arXiv:2102.03077.

\bibitem{YeganehBehroozi2025}
R.~S.~Yeganeh and H.~Behroozi, ``Energy Efficient RSMA-Based LEO Satellite
Communications Assisted by UAV-Mounted BD-Active RIS: A DRL Approach,''
arXiv:2505.04148, 2025.

\bibitem{RiskAwareLEO2024}
``Robust Beamforming for Massive MIMO LEO Satellite Communications: A
Risk-Aware Learning Framework,'' \emph{IEEE} [venue/author details to be
verified before citing in a manuscript --- not independently confirmed],
IEEE Xplore document 10336551.

\bibitem{JiaoFangDing2022}
S.~Jiao, F.~Fang, and Z.~Ding, ``Joint Robust Beamforming Design for
WPT-assisted D2D Communications in MISO-NOMA: Fractional Programming and
Deep Reinforcement Learning,'' arXiv:2209.12253, 2022.

\bibitem{HaEtAl2024}
V.~N.~Ha, J.~C.~Merlano Duncan, E.~Lagunas, J.~Querol, and S.~Chatzinotas,
``Energy-Efficient Precoding and Feeder-Link-Beam Matching Design for
Bent-Pipe SATCOM Systems,'' arXiv:2304.13436, 2023 (revised 2024).

\bibitem{Lohmeyer2016}
W.~Lohmeyer et al., ``Communication Satellite Power Amplifiers: Current and
Future SSPA and TWTA Technologies,'' \emph{International Journal of
Satellite Communications and Networking}, vol.~34, no.~2, pp.~95--113, 2016.

\bibitem{DengEtAl2023}
Deng, Li, Jhaveri, et al., ``Reinforcement-Learning-Based Optimization on
Energy Efficiency in UAV Networks for IoT,'' \emph{IEEE Internet of Things
Journal}, 2023.

\bibitem{SharpeRatioMDP2025}
``Sharpe Ratio Optimization in Markov Decision Processes,'' arXiv:2509.00793,
2025.

\bibitem{Paternain2019}
S.~Paternain, L.~F.~O.~Chamon, M.~Calvo-Fullana, and A.~Ribeiro,
``Constrained Reinforcement Learning Has Zero Duality Gap,'' \emph{Advances
in Neural Information Processing Systems (NeurIPS)}, 2019. arXiv:1910.13393.

\bibitem{DingZhangBasar2020}
D.~Ding, K.~Zhang, and T.~Ba\c{s}ar, ``Natural Policy Gradient Primal-Dual
Method for Constrained Markov Decision Processes,'' \emph{Advances in Neural
Information Processing Systems (NeurIPS)}, 2020. arXiv:2206.02346.

\bibitem{JinTangZhangWang2023}
Jin, Tang, Zhang, and Wang, ``Fractional Deep Reinforcement Learning for
Age-Minimal Mobile Edge Computing,'' arXiv:2312.10418, 2023.

\bibitem{GhadimiLan2013}
S.~Ghadimi and G.~Lan, ``Stochastic First- and Zeroth-Order Methods for
Nonconvex Stochastic Programming,'' \emph{SIAM Journal on Optimization},
vol.~23, no.~4, pp.~2341--2368, 2013.

\bibitem{GhadimiLanZhang2016}
S.~Ghadimi, G.~Lan, and H.~Zhang, ``Mini-Batch Stochastic Approximation
Methods for Nonconvex Stochastic Composite Optimization,'' \emph{Mathematical
Programming}, vol.~155, pp.~267--305, 2016.

\bibitem{StookeAchiamAbbeel2020}
A.~Stooke, J.~Achiam, and P.~Abbeel, ``Responsive Safety in Reinforcement
Learning by PID Lagrangian Methods,'' \emph{International Conference on
Machine Learning (ICML)}, 2020. arXiv:2007.03964.

\end{thebibliography}

\end{document}
